\chapter{The supervised-learning contract}
\label{ch:supervised-contract}

\begin{learningobjectives}
After this chapter, you should be able to:
\begin{itemize}
  \item formulate a supervised-learning task as a statistical and operational contract;
  \item distinguish features, targets, predictions, decisions, and causal claims;
  \item design a split that represents the intended prediction setting; and
  \item construct a baseline and a minimal training-to-feedback path.
\end{itemize}
\end{learningobjectives}

\begin{chapterconnection}
Part I established how data and code cross representation, package, database,
API, and operational boundaries. Part II inserts a learned component into that
system. Learning does not remove those responsibilities. It adds estimation
error, distributional assumptions, adaptive selection, and delayed evidence.
\end{chapterconnection}

\section{Learning begins with a future information boundary}

Let an observation be a pair $(X,Y)$ drawn from a population distribution
$P$. The feature vector $X\in\mathcal{X}$ contains information available at a
specified prediction time. The target $Y\in\mathcal{Y}$ records the outcome to
be estimated. A learning algorithm receives a sample
$D=\{(x_i,y_i)\}_{i=1}^n$ and produces a predictor
$\widehat f_D:\mathcal{X}\rightarrow\mathcal{A}$. The action space
$\mathcal{A}$ may contain numerical estimates, scores, probabilities, or class
labels.

This notation is compact enough to hide the most important design decisions.
What is one observational unit? From which population were units sampled? At
what time is $X$ frozen? How far in the future is $Y$ measured? Which units can
recur? What action consumes the output? The subscript $D$ is a reminder that a
fitted predictor is random: a different sample can produce a different
function.

\begin{definitionbox}{supervised-learning contract}
A supervised-learning contract names the observational unit, target population,
prediction time, target and horizon, permitted information, output semantics,
loss or utility, evaluation design, operational consumer, and outcome-feedback
process. The model class is one field in this contract, not the contract itself.
\end{definitionbox}

Our recurring case concerns battery cells. At the end of an early measurement
window, we may estimate capacity after a later number of cycles, or predict
whether the cell will cross a degradation threshold within a stated horizon.
Measurements collected after the prediction time are forbidden, even if they
appear in the same convenient table.

\section{Risk, empirical risk, and the generalization claim}

For a loss function $L(y,a)$, the population risk of $f$ is
\[
R(f)=\mathbb{E}_{(X,Y)\sim P}[L(Y,f(X))].
\]
The distribution $P$ is generally unknown. Training therefore uses an empirical
quantity such as
\[
\widehat R_D(f)=\frac{1}{n}\sum_{i=1}^n L(y_i,f(x_i)).
\]
Minimizing empirical loss is a computational objective. The scientific claim
is that performance transfers to a named population and use setting. That claim
depends on sampling, dependence, temporal stability, measurement, and the way
the model was selected.

A loss trains or compares functions. A metric communicates some aspect of
behavior. A decision policy converts output into action. These objects can be
different. Logistic loss may fit a probability model, calibration error may
audit its probabilities, and an asymmetric cost rule may choose an intervention
threshold.

\conceptcheck{A laboratory reports excellent mean squared error after randomly
splitting rows, but each physical cell contributes many rows to both subsets.
What population quantity does that experiment estimate, and why might it differ
from performance on a previously unseen cell?}

\section{Prediction is not causal explanation}

A predictor exploits associations that persist between its training and use
settings. It need not represent a mechanism. A variable can improve prediction
because it is a proxy for laboratory, instrument, treatment assignment, or
selection process. Conversely, a causally important variable may add little
predictive information if its variation is small or already encoded elsewhere.

The question ``What will happen?'' differs from ``What would happen if we
intervened?'' Counterfactual conclusions require additional design and
assumptions. High predictive accuracy does not supply them.

\begin{misconceptionbox}
\textbf{Misconception:} a sufficiently accurate model discovers the true
causes. \textbf{Correction:} predictive performance is evidence about a
specified prediction task under a specified evaluation design. Causal claims
require an identification argument that predictive accuracy alone cannot
provide.
\end{misconceptionbox}

\section{Splits simulate the use setting}

A random row split is appropriate only when future uses resemble independent
draws from the same row-level population. Dependence changes the unit of
separation. Repeated measures from one cell call for grouped splitting when the
claim concerns unseen cells. Forecasting calls for training on the past and
evaluating on the future. A new laboratory or manufacturing batch may require a
held-out domain.

Three subsets have different responsibilities. Training data determine fitted
parameters. Validation data guide model and hyperparameter choices. Test data
support a final estimate after those choices stop. Cross-validation rotates
training and validation roles; it does not excuse repeatedly consulting the
test set.

\term{Leakage} occurs when model construction uses information that would not
be legitimately available at prediction time or crosses an evaluation
boundary. Target-derived features are obvious leakage. More subtle cases
include fitting normalization on the full dataset, selecting variables before
cross-validation, using future measurements, and allowing related entities into
both sides of a split.

\begin{warningbox}
Preprocessing learns parameters too. Imputation values, category vocabularies,
scales, feature selection, and representation learning must be fitted using the
training portion of each evaluation fold and then applied to held-out data.
\end{warningbox}

\section{A baseline is an experimental control}

A baseline asks whether learning adds value beyond a simple policy. For
regression, predict the training mean or a physically motivated constant. For
classification, predict class prevalence or the most frequent class, but audit
whether that behavior is useful under imbalance. A time-series baseline may
carry the latest observation forward.

The baseline uses the same split, metric, and information boundary as the
candidate. Otherwise it is not a control. A complex model that fails to beat a
correct baseline has supplied diagnostic evidence: the signal, representation,
sample size, objective, or implementation may not support the claim.

\begin{workedexample}{write the battery prediction contract}
Suppose one row summarizes the first 100 cycles of one cell. The target is
capacity at cycle 300, measured in ampere-hours. The intended use is prediction
for cells from future manufacturing batches immediately after cycle 100.

The feature cutoff excludes all observations after cycle 100. Batch identity
defines an outer temporal split: earlier batches train the system and later
batches evaluate it. All transformations are fitted within the training side.
The baseline predicts the mean training capacity at cycle 300. Mean absolute
error is reported in ampere-hours together with the error distribution by batch.

The contract also states what it does not establish. It does not estimate the
effect of changing a charging protocol, certify safety, or guarantee behavior
for a chemistry absent from the study. Those exclusions are part of the result,
not fine print.
\end{workedexample}

\section{The first vertical slice should be deliberately simple}

Before optimizing an algorithm, prove that the system can reproduce a split,
fit a baseline, serialize an artifact, load it in a separate process, validate
inputs, produce predictions, record model and feature versions, and later join
predictions to outcomes. A dummy model is ideal because unexpected behavior is
easy to detect.

The artifact manifest should identify training data, code revision,
configuration, feature schema, split policy, fitted object, evaluation report,
and creation time. A model file without this context is an opaque binary, not a
reproducible result.

\section{Exercises}

\begin{enumerate}
  \item Write two prediction contracts for the same battery table: one for an
  unseen cell from a known batch and one for a future batch. Explain why their
  splits differ.
  \item Construct three leakage examples: one temporal, one group-based, and
  one caused by preprocessing. For each, identify the invalid information path.
  \item Specify regression and classification baselines for a target with strong
  batch imbalance. State what each baseline controls.
  \item Design an artifact manifest sufficient for another process to reject an
  incompatible feature vector before inference.
\end{enumerate}

\begin{chapterreview}
Supervised learning estimates a rule from labeled observations, but the
defensible object is larger than the fitted rule. The contract names population,
time, target, information, loss, split, consumer, and feedback. Evaluation must
simulate the intended use; baselines act as controls; and the first vertical
slice should validate artifact and evidence flow before model complexity grows.
\end{chapterreview}

\begin{notebooklink}
Lecture 9 notebook 00 turns the contract into an executable schema. Its
companion laboratories introduce leakage-safe splits, baselines, and a minimal
training-to-feedback path for the running battery case.
\end{notebooklink}
